\documentclass[10pt,a4paper]{amsart}
\usepackage[margin=19mm]{geometry}
\usepackage{amsmath,amssymb,mathtools}
\usepackage[scr=rsfs,cal=euler]{mathalfa}
\usepackage{xfrac}
\usepackage[unicode,hidelinks,bookmarks=false]{hyperref}
\newcommand{\ie}{i.e.\ }
\newcommand{\cf}{cf.\ }
\newcommand{\resp}{respectively\ }
\newcommand{\Subsection}{\subsection}
\providecommand{\nobreakdash}{\nobreak\hbox{-}}
\setlength{\emergencystretch}{2em}
\multlinegap0pt
\usepackage{float}
\usepackage{amssymb}
\usepackage{mathtools}
\usepackage[shortlabels]{enumitem}
\newtheorem{lem}{Lemma}
\newcommand{\Cat}[1]{\mathcal{#1}}
\renewcommand{\lim}{\operatorname{lim}}
\newcommand{\op}[1]{#1^{\operatorname{op}}}
\newcommand{\zar}{\operatorname{zar}}
\title{The derived \texorpdfstring{$\infty$}{infinity}-category of Frobenius modules}
\author{Klaus Mattis\, and Timo Weiß}
\date{Source: arXiv:2510.23267v2 (CC BY 4.0)}
\begin{document}
\maketitle

\noindent\emph{Source and licence.} The derived \texorpdfstring{$\infty$}{infinity}-category of Frobenius modules. Version arXiv:2510.23267v2 (CC BY 4.0), distributed by arXiv under the Creative Commons Attribution 4.0 International licence, http://creativecommons.org/licenses/by/4.0/, as recorded in the arXiv licence metadata for this exact version and shown on the abstract page https://arxiv.org/abs/2510.23267v2. Full licence notice: \texttt{licenses/arxiv-2510.23267-CC-BY-4.0.txt}.\par\smallskip

\noindent\textit{Editorial scope.} Counted unit: the article's lem lem:zariski:check-on-affines (source0.tex lines 1278--1285). The article's complete original proof of it is delivered with it (source0.tex lines 1286--1302, one \texttt{proof} environment of 17 lines). No mathematics is written, repaired or re-derived: the statement and the proof are the article's own lines, in the article's own notation, and no retained line is substituted or re-flowed.  No further source-local context is needed: every cross-reference of every retained line resolves inside the two delivered spans.  Nothing was omitted from any retained span, so the package carries no omission record.  No retained line contains a citation, so the wrapper carries no bibliography and no bibliography entry.  No retained line contains a figure environment, an image-inclusion command or drawing code, so no image is delivered and the package declares no asset dependency.  Editorial formatting only, in addition: an AMS \texttt{amsart} preamble that restates the article's own declarations and macros used by the retained lines, namely the environments \texttt{lem} on separate counters, together with the standard packages those lines need.  The retained environments print in this document as Lemma 1 in order of appearance, whereas the article prints the same items with the numbering of its own full document.  The complete native source of the article as distributed in the arXiv e-print for this exact version is bundled unchanged as \texttt{sources/187865/source0.tex}.
\begin{lem} \label{lem:zariski:check-on-affines}
    Let $X$ be a qcqs scheme and let $\Cat C$ be a complete $\infty$-category.
    Let $\Cat F, \Cat G \colon \op X_{\zar} \to \Cat C$ be Zariski sheaves,
    and let $f \colon \Cat F \to \Cat G$ be a morphism.
    Suppose that for every affine open $U \subset X$ 
    the morphism $f_U \colon \Cat F(U) \to \Cat G(U)$ is an equivalence.
    Then $f$ is an equivalence.
\end{lem}

\begin{proof}
    We have to show that for every $V \in X_{\zar}$ the morphism $f_V \colon \Cat F(V) \to \Cat G(V)$ 
    is an equivalence. Let $\Cat U = \{U_i\}_i$ be a finite Zariski cover of $V$, where the $U_i$ are affine.
    
    Assume first that $V$ has affine diagonal. Thus, any intersection of the $U_i$ 
    is also affine.
    In particular, the \v{C}ech nerve $\Cat U^\bullet$ of the cover consists only of affine schemes.
    Since $\Cat F$ is a sheaf, $\Cat F(V) \cong \lim_{\Delta} \Cat F(\Cat U^\bullet)$,
    and similarly for $\Cat G$.
    By assumption, $f$ induces an equivalence
    on the limit diagram and therefore also on the limit.

    Now, let $V$ be an arbitrary object of $X_{\zar}$. Note that any intersection of the $U_i$ 
    is quasi-affine (i.e.\ an open subset of an affine scheme). In particular, it is separated 
    and thus has affine diagonal. Arguing as above, we see that $f$ induces an equivalence 
    on the limit diagrams (using the first part of the proof), and hence on the limit.
\end{proof}

\end{document}
